%! TEX program = xelatex
\documentclass[10pt,compress,aspectratio=169]{beamer}
\usepackage[utf8]{inputenc}
\usepackage[OT1]{fontenc}
\usepackage{amsmath,amsthm,amssymb}
\usepackage{mathpartir}
\usepackage{stmaryrd}
\usepackage{booktabs}
\usepackage{minted}
\usepackage{fontspec}
\setmonofont{JetBrains Mono}

% Load Theme
\usetheme[
    navigation=false,
    frametotal=false,
    maincolor=rptudunkelblau,
    oneColorBoxes,
    signetontitle=true
]{RPTU}

\title{Deriving Interpreters:}
\subtitle{From Big-Step SOS to Abstract Machines}
\date[]{\today}
\author[M.\,Hamido]{Mahmoud M. Hamido}

\institute[]{
Department of Computer Science\\ RPTU Kaiserslautern-Landau\\[0.6ex]

\textbf{Supervised by:}\\
Prof.\ Dr.\ Ralf Hinze\\
Alexander Dinges, M.Sc.
}

\begin{document}
\maketitle

\begin{frame}{Outline}
    \tableofcontents
\end{frame}

\section{Motivation}
\begin{frame}{Motivation: I}
\begin{itemize}
	\item Implementations should faithfully realize language specifications.
	\item A mismatch between specification and implementation is a bug.
	\item Testing only checks finitely many inputs.
	\item Specifications claim behavior for all well-formed programs.
	\item Goal: prove semantic agreement by construction.
\end{itemize}
\end{frame}

\begin{frame}{Motivation: II}
\begin{itemize}
    \item Practice: textual specification + conformance tests.
    \item Limitation: tests cannot prove absence of mismatch.
    \item Approach: encode semantics in Agda, then prove correctness.
    \item Curry--Howard: propositions as types, proofs as programs.
    \item Result: one proof covers all inputs simultaneously.
\end{itemize}
\end{frame}

\begin{frame}{Thesis}
\begin{itemize}
    \item Central claim: derive implementations directly from semantics.
    \item Focus: derive abstract machines from big-step semantics.
    \item Case studies: Arithmetic, Exceptions, and STLC.
    % \item Prove consistency through a chain of equivalences.
\end{itemize}
% \chain
\end{frame}

\section{Background}
\begin{frame}[fragile]{Background: Agda}
\begin{itemize}
    \item Agda is a dependently typed programming language and proof assistant.
    \item Types may depend on values, enabling type-level invariants.
\end{itemize}
\vspace{0.5ex}
\begin{minted}{agda}
data Vec (A : Set) : ℕ → Set where
    []  : Vec A zero
    _∷_ : A → Vec A n → Vec A (suc n)
\end{minted}
\begin{itemize}
    \item Here the index tracks list length in the type.
\end{itemize}
\end{frame}

\begin{frame}[fragile]{Background: Propositions as Types}
\begin{itemize}
    \item Truth and falsehood are represented as data types.
\end{itemize}
\begin{minted}{agda}
data ⊤ : Set where
    tt : ⊤

data ⊥ : Set where

ex-falso : {P : Set} → ⊥ → P
ex-falso ()
\end{minted}
\begin{itemize}
    \item A well-typed term is simultaneously a program and a proof.
\end{itemize}
\end{frame}

\begin{frame}[fragile]{Background: Connectives in Agda}
\begin{minted}{agda}
¬_ : Set → Set
¬ P = P → ⊥

data _×_ (P Q : Set) : Set where
    _,_ : P → Q → P × Q

data Σ (A : Set) (B : A → Set) : Set where
    _,_ : (x : A) → B x → Σ A B

data _⊎_ (P Q : Set) : Set where
    inl : P → P ⊎ Q
    inr : Q → P ⊎ Q
\end{minted}
\begin{itemize}
    \item Negation is implication into falsehood.
    \item Product and sum model conjunction and disjunction.
    \item $\Sigma$-types model dependent pairs (existentials).
\end{itemize}
\end{frame}

\begin{frame}[fragile]{Background: Equality and Naturals}
\begin{minted}{agda}
data _≡_ {A : Set} (x : A) : A → Set where
    refl : x ≡ x

data ℕ : Set where
    zero : ℕ
    suc  : ℕ → ℕ

_+_ : ℕ → ℕ → ℕ
zero  + m = m
suc n + m = suc (n + m)
\end{minted}
\begin{itemize}
    \item Equality is indexed and proved by \texttt{refl}.
    \item Addition is defined structurally on the first argument.
\end{itemize}
\end{frame}

\begin{frame}[fragile]{Background: Helper Lemmas (Rewrite)}
\begin{minted}{agda}
n+0≡n : (n : ℕ) → n + zero ≡ n
n+0≡n zero    = refl
n+0≡n (suc n) rewrite n+0≡n n = refl

n+sm≡n : (n m : ℕ) → n + suc m ≡ suc (n + m)
n+sm≡n zero    m = refl
n+sm≡n (suc n) m rewrite n+sm≡n n m = refl
\end{minted}
\begin{itemize}
    \item \texttt{rewrite} substitutes equals in the current goal.
    \item Each recursive case reduces to \texttt{refl} after rewriting.
\end{itemize}
\end{frame}

\begin{frame}[fragile]{Background: Main Theorem Example}
\begin{minted}{agda}
+-commutative : (n m : ℕ) → n + m ≡ m + n
+-commutative zero    m rewrite n+0≡n m = refl
+-commutative (suc n) m
  rewrite +-commutative n m = sym (n+sm≡n m n)
\end{minted}
\begin{itemize}
    \item Strategy: structural induction on \texttt{n}.
    \item Base case uses \texttt{n+0≡n}.
    \item Step case uses induction hypothesis and helper lemma \texttt{n+sm≡n}.
\end{itemize}
\end{frame}

\begin{frame}[fragile]{Background: Termination Obligation}
\begin{minted}{agda}
diverging : ⊥
diverging = diverging
\end{minted}
\begin{itemize}
    \item Without termination, \texttt{ex-falso} would make logic inconsistent.
    \item Agda accepts only structurally decreasing recursion by default.
    \item \verb|{-# TERMINATING #-}| can bypass this, shifting trust to the programmer.
\end{itemize}
\end{frame}

\begin{frame}{Background: Variables and Semantics}
\begin{itemize}
    \item De Bruijn indices avoid capture and $\alpha$-equivalence issues.
    \item Contexts track variable types: $\Gamma ::= \emptyset \mid \Gamma, T$.
    \item Typing judgment: $\Gamma \vdash e : T$.
    \item Four views of semantics are related in this thesis:
\end{itemize}
% \chain
\end{frame}

\section{Case Study I: Arithmetic}
\begin{frame}[fragile]{Arithmetic: Syntax and Intrinsic Typing}
\begin{minted}{agda}
data Type : Set where
    Nat  : Type
    Bool : Type

data ⊢_ : Type → Set where
    `_            : ⊨ T → ⊢ T
    _⊞_           : ⊢ Nat → ⊢ Nat → ⊢ Nat
    test          : ⊢ Nat → ⊢ Bool
    if_then_else_ : ⊢ Bool → ⊢ T → ⊢ T → ⊢ T
\end{minted}
\begin{itemize}
    \item Ill-typed programs are not representable.
\end{itemize}
\end{frame}

\begin{frame}[fragile]{Arithmetic: Big-Step Semantics}
\begin{minted}{agda}
data _⇓_ : ⊢ T → ⊨ T → Set where
    `_      : (v : ⊨ T) → (` v) ⇓ v
    _⊞_     : e₁ ⇓ n₁ → e₂ ⇓ n₂ → (e₁ ⊞ e₂) ⇓ (n₁ + n₂)
    test    : e ⇓ n → (test e) ⇓ (test' n)
    if-then : e ⇓ true  → e₁ ⇓ v → (if e then e₁ else e₂) ⇓ v
    if-else : e ⇓ false → e₂ ⇓ v → (if e then e₁ else e₂) ⇓ v
\end{minted}

\begin{itemize}
    \item Properties proved: determinism and totality.
\end{itemize}
\end{frame}

\begin{frame}[fragile]{Arithmetic: Deterministic and Total}
\begin{minted}{agda}
deterministic : ∀ {e : ⊢ T} {v₁ v₂ : ⊨ T}
    → e ⇓ v₁ → e ⇓ v₂ → v₁ ≡ v₂

total : ∀ (e : ⊢ T) → ∃ (λ v → e ⇓ v)
\end{minted}
\begin{itemize}
    \item Determinism proof: induction on evaluation derivations.
    \item In conditional cases, contradictory branch choices yield impossible equalities.
    \item Totality proof: structural induction on syntax of expressions.
    \item Recursive calls return both a value and a derivation tree.
\end{itemize}
\end{frame}

\begin{frame}[fragile]{Arithmetic: Reverse Simulation Proof Shape}
\begin{minted}{agda}
simulateᴿ : ∀ {T : Type} {e : ⊢ T} {v : ⊨ T}
  → e ⇾⃰ ` v → e ⇓ v

eq = ⇾⃰-confluence tr' tr ¬step-` ¬step-`
\end{minted}
\begin{itemize}
    \item Start from totality: obtain \texttt{v'} with big-step proof \texttt{e ⇓ v'}.
    \item Use forward simulation to build \texttt{tr' : e ⇾⃰ ` v'}.
    \item Use confluence on \texttt{tr'} and given run \texttt{tr : e ⇾⃰ ` v}.
    \item Since both targets are stuck values, confluence yields \texttt{v' ≡ v}.
    \item Transport the big-step proof across that equality via \texttt{subst}.
\end{itemize}
\end{frame}

\begin{frame}[fragile]{Arithmetic: Small-Step and Simulation}

% todo: how do I fit the proof for this on the slides?
\begin{minted}{agda}
simulate : ∀ {T : Type} {e : ⊢ T} {v : ⊨ T}
    → e ⇓ v → e ⇾⃰ ` v

simulate (` v) = base
simulate (p₁ ⊞ p₂) = ...
simulate (if-then pe p₁) = ...
simulate (if-else pe p₂) = ...
\end{minted}
% \begin{itemize}
%     \item Proof uses lifted congruence ($\xi$) rules over transitive closure.
%     \item The addition case composes three traces: left operand, right operand, then $\beta$-step.
%     \item The if-cases first reduce the guard, then take $\beta$-if-then or $\beta$-if-else.
% \end{itemize}
\end{frame}

\begin{frame}[fragile]{Arithmetic: Denotational Semantics}
\begin{minted}{agda}
⟦_⟧ : ∀ {T} → ⊢ T → ⊨ T
⟦ ` v ⟧                  = v
⟦ e₁ ⊞ e₂ ⟧             = ⟦ e₁ ⟧ + ⟦ e₂ ⟧
⟦ test e ⟧               = test' ⟦ e ⟧
⟦ if e then e₁ else e₂ ⟧ with ⟦ e ⟧
... | true  = ⟦ e₁ ⟧
... | false = ⟦ e₂ ⟧
\end{minted}
\begin{itemize}
    \item Equivalence theorem: $e \Downarrow v \leftrightarrow \llbracket e \rrbracket \equiv v$.
\end{itemize}
\end{frame}

\begin{frame}[fragile]{Arithmetic: Abstract Machine Structure}
\begin{minted}{agda}
data Frame : Type → Type → Set where
    _𝚎𝚕𝚜𝚎_ : ⊢ T → ⊢ T → Frame T Bool
    _⊞₀_   : Hole → ⊢ Nat → Frame Nat Nat
    _⊞₁_   : ⊨ Nat → Hole → Frame Nat Nat
    test   : Hole → Frame Bool Nat

data State (Answer : Type) : Set where
    ⊢_↓_ : Stack Answer T → ⊨ T → State Answer
    ⊢_↑_ : Stack Answer T → ⊢ T → State Answer
\end{minted}
\end{frame}

\begin{frame}[fragile]{Arithmetic: Completeness Proof (Machine)}
\begin{minted}{agda}
complete : ∀ (e : ⊢ T) (v : ⊨ T) (stack : Stack Answer T)
    → e ⇓ v → (⊢ stack ↑ e) ⇾⃰ (⊢ stack ↓ v)
\end{minted}
\begin{itemize}
    \item Proof by induction on the big-step derivation \texttt{e ⇓ v}.
    \item Generalized over arbitrary \texttt{stack}, not only \texttt{[]}, because recursive calls push frames.
    \item Each big-step constructor is matched by a short machine trace.
    \item Composition uses transitivity of \texttt{⇾⃰} to concatenate subtraces.
\end{itemize}
\end{frame}

\begin{frame}[fragile]{Arithmetic: Correctness Proof (Machine)}
\begin{minted}{agda}
correct : ∀ (e : ⊢ T) (v : ⊨ T)
    → (⊢ [] ↑ e) ⇾⃰ (⊢ [] ↓ v) → e ⇓ v
\end{minted}
\begin{itemize}
    \item Start from totality of big-step to obtain some \texttt{v'} and proof \texttt{e ⇓ v'}.
    \item By completeness, build a canonical machine run \texttt{⊢ [] ↑ e ⇾⃰ ⊢ [] ↓ v'}.
    \item Compare canonical run with given run \texttt{⊢ [] ↑ e ⇾⃰ ⊢ [] ↓ v} using confluence.
    \item Project state equality to value equality and transport \texttt{e ⇓ v'} to \texttt{e ⇓ v}.
    \item Equality of final states is projected to equality of returned values.
\end{itemize}
\end{frame}

\begin{frame}[fragile]{Arithmetic: Example Completeness Case (if)}
\begin{minted}{agda}
complete (if e then e₁ else e₂) v stack (if-then p p₁) =
  step step-𝚒𝚏
  (⇾⃰-transitive (complete e _ (stack ∷ (e₁ 𝚎𝚕𝚜𝚎 e₂)) p)
  (step step-𝚝𝚑𝚎𝚗
  (complete e₁ v stack p₁)))
\end{minted}
\begin{itemize}
    \item Push branch frame, evaluate guard under extended stack.
    \item After guard returns \texttt{true}, take \texttt{step-then}.
    \item Continue by induction hypothesis on selected branch.
\end{itemize}
\end{frame}

\section{Case Study II: Exceptions}
\begin{frame}[fragile]{Exceptions: Syntax Extension}
\begin{minted}{agda}
data Exception : Set where
    err : Exception

data ⊢_ where
    ...
    throw      : Exception → ⊢ T
    try_catch_ : ⊢ T → (Exception → ⊢ T) → ⊢ T
\end{minted}
\begin{itemize}
    \item Evaluation now yields either a value or an exception.
\end{itemize}
\end{frame}

\begin{frame}[fragile]{Exceptions: Big-Step with Results}
\begin{minted}{agda}
data Result (T : Type) : Set where
    return : ⊨ T → Result T
    raise  : Exception → Result T

data _⇓ʳ_ : ⊢ T → Result T → Set where
    throw-rule : throw ex ⇓ʳ raise ex
    try-return : e ⇓ v → (try e catch h) ⇓ v
    try-catch  : e ⇑ ex → h ex ⇓ʳ r → (try e catch h) ⇓ʳ r
\end{minted}
\begin{itemize}
    \item Every construct has success and propagation behavior.
    \item \texttt{try-return}: handler is skipped when body succeeds.
    \item \texttt{try-catch}: body exception triggers handler evaluation.
\end{itemize}
\end{frame}

\begin{frame}[fragile]{Exceptions: Small-Step Bridge}
\begin{minted}{agda}
resultExpr : ∀ {T : Type} → Result T → ⊢ T
resultExpr (return v) = ` v
resultExpr (raise ex) = throw ex

simulate : ∀ {T : Type} {e : ⊢ T} {r : Result T}
    → e ⇓ʳ r → e ⇾⃰ resultExpr r
\end{minted}
\begin{itemize}
    \item Final expressions are stuck: neither quoted values nor throws step further.
    \item This stuckness lemma is required for the reverse direction via confluence.
    \item Lifted \texttt{ξ-try} composes body reduction with outer try-context.
\end{itemize}
\end{frame}

\begin{frame}[fragile]{Exceptions: Abstract Machine with Unwind Mode}
\begin{minted}{agda}
data State (Answer : Type) : Set where
    ⊢_↑_ : Stack Answer T → ⊢ T      → State Answer
    ⊢_↓_ : Stack Answer T → ⊨ T      → State Answer
    ⊢_☇_ : Stack Answer T → Exception → State Answer

step-throw     : ⊢ stack ↑ throw ex ⇾ ⊢ stack ☇ ex
step-try       : ⊢ stk ↑ (try e catch h) ⇾ ⊢ stk ∷ catch h ↑ e
step-catch-exn : ⊢ stack ∷ catch h ☇ ex ⇾ ⊢ stack ↑ h ex
\end{minted}
\end{frame}

\begin{frame}[fragile]{Exceptions: Completeness Proofs}
\begin{minted}{agda}
complete     : e ⇓ v  → (⊢ stk ↑ e) ⇾⃰ (⊢ stk ↓ v)
complete-exn : e ⇑ ex → (⊢ stk ↑ e) ⇾⃰ (⊢ stk ☇ ex)
\end{minted}
\begin{itemize}
    \item Two mutually recursive lemmas are required: value case and exception case.
    \item \texttt{try-catch} connects both directions:
    \item \texttt{complete} calls \texttt{complete-exn} when the body raises.
    \item \texttt{complete-exn} calls \texttt{complete} when evaluating the handler.
    \item Base outcomes differ: values end in \texttt{↓}, exceptions end in \texttt{☇}.
\end{itemize}
\end{frame}

\begin{frame}[fragile]{Exceptions: Correctness Proofs}
\begin{minted}{agda}
correct     : (⊢ [] ↑ e) ⇾⃰ (⊢ [] ↓ v)  → e ⇓ v
correct-exn : (⊢ [] ↑ e) ⇾⃰ (⊢ [] ☇ ex) → e ⇑ ex
\end{minted}
\begin{itemize}
    \item Proof begins with totality: \texttt{e} evaluates to either \texttt{return v'} or \texttt{raise ex'}.
    \item In the matching terminal mode, confluence yields equality of terminal payloads.
    \item In the mismatching mode (\texttt{↓} vs \texttt{☇}), confluence yields constructor contradiction.
    \item This contradiction is discharged with \texttt{ex-falso}, giving the desired result.
\end{itemize}
\end{frame}

\begin{frame}[fragile]{Exceptions: Correctness Case Analysis}
\begin{minted}{agda}
correct e v tr with total e
... | return v' ﹐ p = ... confluence ...
... | raise ex  ﹐ p = ... ex-falso (↓≢☇ ...)
\end{minted}
\begin{itemize}
    \item If totality yields \texttt{return v'}, canonical completeness run ends in \texttt{↓ v'}.
    \item Confluence with given run to \texttt{↓ v} gives \texttt{v' = v}.
    \item If totality yields \texttt{raise ex}, completeness-exn ends in \texttt{☇ ex}.
    \item Confluence with a \texttt{↓}-run yields constructor clash, discharged by absurdity.
\end{itemize}
\end{frame}

\section{Case Study III: STLC}
\begin{frame}[fragile]{STLC: Syntax and De Bruijn Membership}
\begin{minted}{agda}
data _⊢_ (Γ : Context) : Type → Set where
    Nat  : ℕ → Γ ⊢ Nat
    Bool : 𝔹 → Γ ⊢ Bool
    Var  : T ∈ Γ → Γ ⊢ T
    ƛ_   : (Γ , T₁) ⊢ T₂ → Γ ⊢ (T₁ ⇒ T₂)
    _·_  : Γ ⊢ (T₁ ⇒ T₂) → Γ ⊢ T₁ → Γ ⊢ T₂
\end{minted}
\begin{itemize}
    \item Context-indexed syntax guarantees scope and typing discipline.
\end{itemize}
\end{frame}

\begin{frame}[fragile]{STLC: Closures and Big-Step Application}
\begin{minted}{agda}
record Closure (A B : Type) : Set where
    constructor ⟨_,_⟩
    field
        {Γᶜ} : Context
        δᶜ   : Environment Γᶜ
        eᶜ   : (Γᶜ , A) ⊢ B

app : δ ⊢ e₀ ⇓ ⟨ δᶜ , eᶜ ⟩ → δ ⊢ e₁ ⇓ v₁
     → (δᶜ , v₁) ⊢ eᶜ ⇓ v → δ ⊢ e₀ · e₁ ⇓ v
\end{minted}
\begin{itemize}
    \item Application first evaluates function position to a closure.
    \item Then evaluates the argument and runs the body in captured environment.
\end{itemize}
\end{frame}

\begin{frame}[fragile]{STLC: Totality and Termination Note}
\begin{minted}{agda}
{-# TERMINATING #-}
total : ∀ (e : Γ ⊢ T) → ∃ (λ v → δ ⊢ e ⇓ v)

total (Var x) = ...
total (ƛ e)   = ...
total (e₁ · e₂) with total e₁ | total e₂
... | ... = ...
\end{minted}
\begin{itemize}
    \item STLC is strongly normalizing, but this development uses a termination pragma.
    \item Reason: recursive call proceeds through closure body, not obviously structural to Agda.
    \item So accepted on programmer trust rather than checker's structural criterion.
\end{itemize}
\end{frame}

\begin{frame}[fragile]{STLC: Abstract Machine for Application}
\begin{minted}{agda}
data Frame : Type → Type → Set where
    app₁ : Hole → Γ ⊢ A      → Frame B (A ⇒ B)
    app₂ : ⊨ (A ⇒ B) → Hole → Frame B A
    app₃ : Environment Γ → ⊨ (A ⇒ B) → ⊨ A → Hole → Frame B B

step-·  : (δ ⊢ stack ↑ e₁ · e₂) ⇾ (δ ⊢ stack ∷ app₁ ◌ e₂ ↑ e₁)
step-·₂ : ... ⇾ ((δᶜ , v₂) ⊢ stack ∷ app₃ δ ⟨ δᶜ , eᶜ ⟩ v₂ ◌ ↑ eᶜ)
step-·₃ : ... ⇾ (δ′ ⊢ stack ↓ v₃)
\end{minted}
\end{frame}

\begin{frame}[fragile]{STLC: Proof Status in This Thesis}
\begin{minted}{agda}
complete : (δ ⊢ stack ↑ e ⇾⃰  δ ⊢ stack ↓ v)
\end{minted}
\begin{itemize}
    \item Completed: big-step semantics, denotational semantics, and machine correspondence.
    \item Completed: machine completeness/correctness follows arithmetic-style confluence argument.
    \item Open: big-step $\leftrightarrow$ small-step proof is not finished in this thesis.
    \item Main obstacle: bridge between substitution-based small-step and environment-based big-step.
\end{itemize}
\end{frame}

\begin{frame}[fragile]{STLC: Why Small-Step Equivalence Remains Open}
\begin{minted}{agda}
embed : ∀ {T : Type} → ⊨ T → ∅ ⊢ T
embed {T = T₁ ⇒ T₂} ⟨ δᶜ , eᶜ ⟩ = ƛ ?
\end{minted}
\begin{itemize}
    \item To state simulation, closure values must be re-embedded as syntax.
    \item The captured environment in closures prevents a direct embedding.
    \item Small-step uses substitution; big-step uses environment lookup.
    \item Proving these two mechanisms equivalent is left as future work.
\end{itemize}
\end{frame}

\section{Related Work}
\begin{frame}{Related Work}
\begin{itemize}
    \item Hutton derives machines by CPS transformation and defunctionalization.
    \item This thesis derives machines from big-step semantics by reifying frames.
    \item Conceptual correspondence: frame $\approx$ defunctionalized continuation.
    \item Hutton and Wright also study exceptions and verified compilation.
\end{itemize}
\end{frame}

\section{Conclusion}
\begin{frame}{Conclusion: Summary}
\begin{itemize}
    \item A systematic derivation route from semantics to abstract machines.
    \item Demonstrated for arithmetic, exceptions, and STLC.
    \item Established equivalence links across semantic styles.
    \item Proved machine consistency with big-step semantics in each completed case study.
\end{itemize}
% \chain
\end{frame}

\begin{frame}{Conclusion: Limitations and Future Work}
\begin{itemize}
    \item Current route assumes totality and does not model divergence directly.
    \item STLC totality currently uses a trusted termination pragma.
    \item STLC small-step $\leftrightarrow$ big-step equivalence is open.
    \item Future directions:
    \begin{itemize}
        \item Coinductive big-step or small-step-first derivation for divergence.
        \item Strengthen STLC termination and equivalence proofs.
        \item Extend to richer effects and language features.
    \end{itemize}
\end{itemize}
\end{frame}

\end{document}
